
\documentclass{article}
\usepackage{colortbl}
\usepackage{makecell}
\usepackage{multirow}
\usepackage{supertabular}

\begin{document}

\newcounter{utterance}

\centering \large Interaction Transcript for game `wordle', experiment `high\_frequency\_words\_no\_clue\_no\_critic', episode 2 with nemotron{-}3{-}ultra.
\vspace{24pt}

{ \footnotesize  \setcounter{utterance}{1}
\setlength{\tabcolsep}{0pt}
\begin{supertabular}{c@{$\;$}|p{.15\linewidth}@{}p{.15\linewidth}p{.15\linewidth}p{.15\linewidth}p{.15\linewidth}p{.15\linewidth}}
    \# & \multicolumn{2}{c}{Player} && \multicolumn{2}{c}{Game Master} \\
    \hline

    \theutterance \stepcounter{utterance}  
    & & & \multicolumn{4}{p{0.6\linewidth}}{
        \cellcolor[rgb]{0.9,0.9,0.9}{
            \makecell[{{p{\linewidth}}}]{
                \texttt{\tiny{[P1$\langle$GM]}}
                \texttt{Вы — языковой волшебник, который любит угадывать слова, используя заданные правила.} \\
\\ 
\texttt{Добро пожаловать в Wordle! У вас есть шесть попыток, чтобы угадать целевое слово — допустимое слово из пяти букв. Пожалуйста, используйте теги "объяснение:" и "догадка:", чтобы дать краткое объяснение для каждой попытки.} \\
\\ 
\texttt{Например, если ваша догадка — "слово", ваш ответ должен быть} \\
\texttt{объяснение: это распространённое слово из пяти букв, и я начинаю угадывать с этого слова.} \\
\texttt{догадка: слово} \\
\\ 
\texttt{После каждой попытки ваш ответ будет проверен, и вы получите обратную связь, показывающую, какие буквы верны (зелёный), какие буквы верны, но стоят не на своём месте (жёлтый), и какие буквы неверны (красный). Эта обратная связь может быть полезна для определения того, какие буквы включить или исключить в следующей попытке.} \\
\\ 
\texttt{Например, обратная связь для "слово" может быть:} \\
\texttt{оценка\_догадки: с<жёлтый> л<жёлтый> о<зелёный> в<жёлтый> о<красный>} \\
\\ 
\texttt{В explanation должны быть детали о том, как оценка\_догадки используется, чтобы прийти к новой догадке.} \\
\\ 
\texttt{Не генерируйте никакого другого текста.} \\
\\ 
\texttt{Начнём с вашей первой догадки.} \\
            }
        }
    }
     \\ \\

    \theutterance \stepcounter{utterance}  
    & \multicolumn{4}{p{0.6\linewidth}}{
        \cellcolor[rgb]{0.561,0.737,0.561}{
            \makecell[{{p{\linewidth}}}]{
                \texttt{\tiny{[P1$\rangle$GM]}}
                \texttt{объяснение: Я начинаю со слова "стоит", так как оно содержит часто встречающиеся в русском языке буквы: С, Т, О, И. Это хорошее стартовое слово для получения максимальной информации.} \\
\texttt{догадка: стоит} \\
            }
        }
    }
    & & \\ \\

    \theutterance \stepcounter{utterance}  
    & & & \multicolumn{2}{p{0.3\linewidth}}{
        \cellcolor[rgb]{0.9,0.9,0.9}{
            \makecell[{{p{\linewidth}}}]{
                \texttt{\tiny{[GM$|$GM]}}
                \texttt{attempts = 1} \\
\texttt{target\_word = сосед} \\
\texttt{guess = стоит} \\
\texttt{guess\_feedback = с<green> т<red> о<yellow> и<red> т<red>} \\
            }
        }
    }
    & & \\ \\

    \theutterance \stepcounter{utterance}  
    & & & \multicolumn{4}{p{0.6\linewidth}}{
        \cellcolor[rgb]{0.9,0.9,0.9}{
            \makecell[{{p{\linewidth}}}]{
                \texttt{\tiny{[P1$\langle$GM]}}
                \texttt{оценка\_догадки: с<зелёный> т<красный> о<жёлтый> и<красный> т<красный>} \\
\\ 
\texttt{Дайте ответ только в этом формате.} \\
\texttt{объяснение: подробности} \\
\texttt{догадка: слово} \\
            }
        }
    }
     \\ \\

    \theutterance \stepcounter{utterance}  
    & \multicolumn{4}{p{0.6\linewidth}}{
        \cellcolor[rgb]{0.561,0.737,0.561}{
            \makecell[{{p{\linewidth}}}]{
                \texttt{\tiny{[P1$\rangle$GM]}}
                \texttt{объяснение: Буква 'с' на первом месте верна (зелёный). Буква 'о' есть в слове, но не на второй позиции (жёлтый). Буквы 'т' и 'и' отсутствуют в слове (красный). Поскольку 'т' красная и стоит на 2{-}й и 5{-}й позициях, их исключаем. 'О' не может быть на 2{-}й позиции, пробуем поставить её на 3{-}ю. Подбираю слово на 'с', с 'о' на 3{-}й позиции, без 'т' и 'и' — "склон".} \\
\texttt{догадка: склон} \\
            }
        }
    }
    & & \\ \\

    \theutterance \stepcounter{utterance}  
    & & & \multicolumn{2}{p{0.3\linewidth}}{
        \cellcolor[rgb]{0.9,0.9,0.9}{
            \makecell[{{p{\linewidth}}}]{
                \texttt{\tiny{[GM$|$GM]}}
                \texttt{attempts = 2} \\
\texttt{target\_word = сосед} \\
\texttt{guess = склон} \\
\texttt{guess\_feedback = с<green> к<red> л<red> о<yellow> н<red>} \\
            }
        }
    }
    & & \\ \\

    \theutterance \stepcounter{utterance}  
    & & & \multicolumn{4}{p{0.6\linewidth}}{
        \cellcolor[rgb]{0.9,0.9,0.9}{
            \makecell[{{p{\linewidth}}}]{
                \texttt{\tiny{[P1$\langle$GM]}}
                \texttt{оценка\_догадки: с<зелёный> к<красный> л<красный> о<жёлтый> н<красный>} \\
\\ 
\texttt{Дайте ответ только в этом формате.} \\
\texttt{объяснение: подробности} \\
\texttt{догадка: слово} \\
            }
        }
    }
     \\ \\

    \theutterance \stepcounter{utterance}  
    & \multicolumn{4}{p{0.6\linewidth}}{
        \cellcolor[rgb]{0.561,0.737,0.561}{
            \makecell[{{p{\linewidth}}}]{
                \texttt{\tiny{[P1$\rangle$GM]}}
                \texttt{объяснение: Буква 'с' на первом месте верна. Буква 'о' есть в слове, но не на 2{-}й и не на 4{-}й позиции (была жёлтой в обоих случаях), значит она может быть только на 3{-}й или 5{-}й. Буквы 'к', 'л', 'н', 'т', 'и' исключены. Пробую слово с 'о' на 3{-}й позиции и частыми буквами 'п', 'р', 'а' — "спорт".} \\
\texttt{догадка: спорт} \\
            }
        }
    }
    & & \\ \\

    \theutterance \stepcounter{utterance}  
    & & & \multicolumn{2}{p{0.3\linewidth}}{
        \cellcolor[rgb]{0.9,0.9,0.9}{
            \makecell[{{p{\linewidth}}}]{
                \texttt{\tiny{[GM$|$GM]}}
                \texttt{attempts = 3} \\
\texttt{target\_word = сосед} \\
\texttt{guess = спорт} \\
\texttt{guess\_feedback = с<green> п<red> о<yellow> р<red> т<red>} \\
            }
        }
    }
    & & \\ \\

    \theutterance \stepcounter{utterance}  
    & & & \multicolumn{4}{p{0.6\linewidth}}{
        \cellcolor[rgb]{0.9,0.9,0.9}{
            \makecell[{{p{\linewidth}}}]{
                \texttt{\tiny{[P1$\langle$GM]}}
                \texttt{оценка\_догадки: с<зелёный> п<красный> о<жёлтый> р<красный> т<красный>} \\
\\ 
\texttt{Дайте ответ только в этом формате.} \\
\texttt{объяснение: подробности} \\
\texttt{догадка: слово} \\
            }
        }
    }
     \\ \\

    \theutterance \stepcounter{utterance}  
    & \multicolumn{4}{p{0.6\linewidth}}{
        \cellcolor[rgb]{0.561,0.737,0.561}{
            \makecell[{{p{\linewidth}}}]{
                \texttt{\tiny{[P1$\rangle$GM]}}
                \texttt{объяснение: Буква 'с' на первом месте верна. Буква 'о' жёлтая на 3{-}й позиции, значит она не на 3{-}й. Так как ранее она была жёлтой на 2{-}й и 4{-}й, единственное оставшееся место для 'о' — 5{-}я позиция. Буквы 'п', 'р', 'т', 'к', 'л', 'н', 'и' исключены. Пробую слово на 'с', оканчивающееся на 'о', с допустимыми буквами 'в', 'а', 'х' — "свахо".} \\
\texttt{догадка: свахо} \\
            }
        }
    }
    & & \\ \\

    \theutterance \stepcounter{utterance}  
    & & & \multicolumn{2}{p{0.3\linewidth}}{
        \cellcolor[rgb]{0.9,0.9,0.9}{
            \makecell[{{p{\linewidth}}}]{
                \texttt{\tiny{[GM$|$GM]}}
                \texttt{Error: The guessed word is not a valid word for this game.} \\
            }
        }
    }
    & & \\ \\

    \theutterance \stepcounter{utterance}  
    & & & \multicolumn{4}{p{0.6\linewidth}}{
        \cellcolor[rgb]{0.9,0.9,0.9}{
            \makecell[{{p{\linewidth}}}]{
                \texttt{\tiny{[P1$\langle$GM]}}
                \texttt{Ваша догадка не является допустимым словом для этой игры. Пожалуйста, попробуйте ещё раз.} \\
\\ 
\texttt{Дайте ответ только в этом формате.} \\
\texttt{объяснение: подробности} \\
\texttt{догадка: слово} \\
            }
        }
    }
     \\ \\

    \theutterance \stepcounter{utterance}  
    & \multicolumn{4}{p{0.6\linewidth}}{
        \cellcolor[rgb]{0.561,0.737,0.561}{
            \makecell[{{p{\linewidth}}}]{
                \texttt{\tiny{[P1$\rangle$GM]}}
                \texttt{объяснение: Слово "свахо" не принято. Буква 'о' должна быть на 5{-}м месте (зелёная), 'с' на 1{-}м (зелёная). Исключены: т, и, к, л, н, п, р. Пробую другое слово на 'с...о' с допустимыми буквами: "сабо" (не 5 букв), "саво" (нет). Выбираю существительное "сага" — нет. Пробую "саксо" — 'к' исключена. Правильное слово с буквами с, а, б, о: "сабо" — 4 буквы. Беру слово "сапог" — 'п' исключена. Остаётся "сахар" — 'р' исключена. Слово "сапёр" — 'р' исключена. Слово "саженец" — долго. Выбираю "самую" — 'м', 'ю' новые. Но лучше "салама" — нет. Подбираю реальное слово: "саван" (нет), "савуа" (нет). Пробую "сага" — нет. Беру "сабля" — 'л', 'я' новые, 'б' новая. Слово "сабля" существует.} \\
\texttt{догадка: сабля} \\
            }
        }
    }
    & & \\ \\

    \theutterance \stepcounter{utterance}  
    & & & \multicolumn{2}{p{0.3\linewidth}}{
        \cellcolor[rgb]{0.9,0.9,0.9}{
            \makecell[{{p{\linewidth}}}]{
                \texttt{\tiny{[GM$|$GM]}}
                \texttt{attempts = 4} \\
\texttt{target\_word = сосед} \\
\texttt{guess = сабля} \\
\texttt{guess\_feedback = с<green> а<red> б<red> л<red> я<red>} \\
            }
        }
    }
    & & \\ \\

    \theutterance \stepcounter{utterance}  
    & & & \multicolumn{4}{p{0.6\linewidth}}{
        \cellcolor[rgb]{0.9,0.9,0.9}{
            \makecell[{{p{\linewidth}}}]{
                \texttt{\tiny{[P1$\langle$GM]}}
                \texttt{оценка\_догадки: с<зелёный> а<красный> б<красный> л<красный> я<красный>} \\
\\ 
\texttt{Дайте ответ только в этом формате.} \\
\texttt{объяснение: подробности} \\
\texttt{догадка: слово} \\
            }
        }
    }
     \\ \\

    \theutterance \stepcounter{utterance}  
    & \multicolumn{4}{p{0.6\linewidth}}{
        \cellcolor[rgb]{0.561,0.737,0.561}{
            \makecell[{{p{\linewidth}}}]{
                \texttt{\tiny{[P1$\rangle$GM]}}
                \texttt{объяснение: Буква 'с' на 1{-}м месте верна, 'о' на 5{-}м месте верна (вывод из прошлых ходов). Буквы 'а', 'б', 'л', 'я' (из "сабля"), а также 'т', 'и', 'к', 'н', 'п', 'р' (из ранних ходов) исключены. Остаются гласные 'у', 'ы', 'э', 'ё', 'ю' и согласные 'в', 'г', 'д', 'ж', 'з', 'м', 'с', 'х', 'ц', 'ч', 'ш', 'щ', 'ф'. Пробую слово "сухо" (4 буквы) {-}> "сухой" (5 букв, оканчивается на 'й', а не 'о'). Пробую "судно" — 'д', 'н' ('н' исключена). Пробую "суета" — 'т', 'а' исключены. Пробую "сумас" — не оканчивается на 'о'. Правильная структура: С \_ \_ \_ О. Пробую "свежо" (краткое прилагательное, 5 букв).} \\
\texttt{догадка: свежо} \\
            }
        }
    }
    & & \\ \\

    \theutterance \stepcounter{utterance}  
    & & & \multicolumn{2}{p{0.3\linewidth}}{
        \cellcolor[rgb]{0.9,0.9,0.9}{
            \makecell[{{p{\linewidth}}}]{
                \texttt{\tiny{[GM$|$GM]}}
                \texttt{attempts = 5} \\
\texttt{target\_word = сосед} \\
\texttt{guess = свежо} \\
\texttt{guess\_feedback = с<green> в<red> е<yellow> ж<red> о<yellow>} \\
            }
        }
    }
    & & \\ \\

    \theutterance \stepcounter{utterance}  
    & & & \multicolumn{4}{p{0.6\linewidth}}{
        \cellcolor[rgb]{0.9,0.9,0.9}{
            \makecell[{{p{\linewidth}}}]{
                \texttt{\tiny{[P1$\langle$GM]}}
                \texttt{оценка\_догадки: с<зелёный> в<красный> е<жёлтый> ж<красный> о<жёлтый>} \\
\\ 
\texttt{Дайте ответ только в этом формате.} \\
\texttt{объяснение: подробности} \\
\texttt{догадка: слово} \\
            }
        }
    }
     \\ \\

    \theutterance \stepcounter{utterance}  
    & \multicolumn{4}{p{0.6\linewidth}}{
        \cellcolor[rgb]{0.561,0.737,0.561}{
            \makecell[{{p{\linewidth}}}]{
                \texttt{\tiny{[P1$\rangle$GM]}}
                \texttt{объяснение: Буква 'с' на 1{-}м месте верна. Буква 'е' есть в слове, но не на 3{-}й позиции (жёлтый). Буква 'о' есть в слове, но не на 5{-}й позиции (жёлтый) — это критически важно: значит, 'о' стоит не на 5{-}м месте. Ранее 'о' была жёлтой на 2{-}й, 3{-}й и 4{-}й позициях. Единственное оставшееся место для 'о' — 1{-}я позиция, но она занята 'с'. Противоречие указывает на то, что в слове ДВЕ буквы 'о'. Одна 'о' зафиксирована на 5{-}м месте (зелёная в будущем), вторая 'о' "плавает" (жёлтые на 2,3,4). Буквы 'в', 'ж' исключены. Пробую слово с двумя 'о': "сОООо" нет. "сОлОо" — 'л' исключена. "сОрОо" — 'р' исключена. "сОкОо" — 'к' исключена. "сОмОо" — 'м' новая. Слово "сомое" не существует. Пробую "соёмо" — 'ё', 'м'. Но 'о' на 5{-}м месте должна быть зелёной, а здесь жёлтая. Значит, в слове НЕТ 'о' на 5{-}м месте. Пересматриваю логику: если 'о' жёлтая на 5{-}м месте, её нет на 5{-}м. Где она? На 2,3,4 — жёлтые. На 1{-}й — 'с'. Значит, 'о' в слове НЕТ? Но она жёлтая. Значит, в слове две 'о': одна на 2/3/4, вторая тоже на 2/3/4. Пробую "сОООк" — нет. "сОкОл" — 'к','л' нет. "сОмОн" — 'н' нет. "сОвОк"} \\
            }
        }
    }
    & & \\ \\

    \theutterance \stepcounter{utterance}  
    & & & \multicolumn{2}{p{0.3\linewidth}}{
        \cellcolor[rgb]{0.9,0.9,0.9}{
            \makecell[{{p{\linewidth}}}]{
                \texttt{\tiny{[GM$|$GM]}}
                \texttt{Error: The response should contain the 'догадка:' keyword exactly once.} \\
            }
        }
    }
    & & \\ \\

    \theutterance \stepcounter{utterance}  
    & & & \multicolumn{2}{p{0.3\linewidth}}{
        \cellcolor[rgb]{0.9,0.9,0.9}{
            \makecell[{{p{\linewidth}}}]{
                \texttt{\tiny{[GM$|$GM]}}
                \texttt{game\_result = ABORT} \\
            }
        }
    }
    & & \\ \\

\end{supertabular}
}

\end{document}
